\documentclass[lualatex,ja=standard]{bxjsarticle}
\usepackage{amsmath,amssymb}
\usepackage{enumitem}
\setlist[enumerate]{itemsep=0.6em, topsep=0.5em}
\pagestyle{empty}

\begin{document}

\noindent\textbf{数IA 計算演習（少し難しめ） 解答}
\hrule
\vspace{0.8em}

\begin{enumerate}
\item
\[
(2x-3)^2-(x+1)(x-5)
=4x^2-12x+9-(x^2-4x-5)
=3x^2-8x+14
\]

\item 分母を払うと
\[
(x+2)^2-(2x-1)(x-1)=3
\]
\[
x^2+4x+4-(2x^2-3x+1)=3
\]
\[
-x^2+7x+3=3
\Rightarrow -x^2+7x=0
\Rightarrow x(x-7)=0
\]
\[
x=0,\ 7
\]
ただし $x\neq 1,-2$ より、どちらも適する。

\item
\[
3x^2-2x-1=(3x+1)(x-1)=0
\]
\[
x=1,\ -\frac13
\]

\item 重解条件より判別式 $D=0$。
\[
(-4a)^2-4\cdot 1\cdot(4a+5)=0
\]
\[
16a^2-16a-20=0
\Rightarrow 4a^2-4a-5=0
\]
\[
a=\frac{4\pm\sqrt{16+80}}{8}=\frac{4\pm 4\sqrt6}{8}=\frac{1\pm\sqrt6}{2}
\]

\item
\[
y=x^2-6x+11=(x-3)^2+2
\]
頂点は $x=3$（区間内）なので最小値は $2$。
端点を確認すると
\[
y(1)=6,\quad y(5)=6
\]
よって最大値は $6$（$x=1,5$）。

\item 第2象限なので $\sin\theta>0$。
\[
\sin\theta=\sqrt{1-\cos^2\theta}=\sqrt{1-\frac{25}{169}}=\frac{12}{13}
\]
\[
\tan\theta=\frac{\sin\theta}{\cos\theta}=\frac{12/13}{-5/13}=-\frac{12}{5}
\]

\item 余弦定理より
\[
BC^2=AB^2+AC^2-2\cdot AB\cdot AC\cos A
\]
\[
=7^2+5^2-2\cdot 7\cdot 5\cdot\cos120^\circ
=49+25-70\left(-\frac12\right)
=109
\]
\[
BC=\sqrt{109}
\]

\item
\[
{}_8P_3=8\cdot 7\cdot 6=336,\quad {}_8C_3=\frac{8\cdot 7\cdot 6}{3\cdot 2\cdot 1}=56
\]
\[
{}_8P_3-{}_8C_3=336-56=280
\]

\item 全体は
\[
\binom{9}{3}=84
\]
青玉を含まない（赤4個・白3個のみから3個）場合は
\[
\binom{7}{3}=35
\]
したがって求める通り数は
\[
84-35=49
\]

\item 全事象は $36$ 通り。
和が9以上は
\[
(3,6),(4,5),(5,4),(6,3),
(4,6),(5,5),(6,4),
(5,6),(6,5),
(6,6)
\]
の10通り。
このうち積が奇数なのは両方奇数のときのみで、該当は
\[
(5,5)
\]
の1通り。
よって条件を満たすのは $10-1=9$ 通り。
\[
\frac{9}{36}=\frac14
\]
\end{enumerate}

\end{document}
